\documentclass[11pt, oneside]{article} 
\usepackage{geometry}
\geometry{letterpaper} 
\usepackage{graphicx}
	
\usepackage{amssymb}
\usepackage{amsmath}
\usepackage{parskip}
\usepackage{color}
\usepackage{hyperref}

\graphicspath{{/Users/telliott/Dropbox/Github-Math/figures/}}

\title{Fundamental similarity}
\date{}

\begin{document}
\maketitle
\Large

%[my-super-duper-separator]

\subsection*{problem}

\begin{center} \includegraphics [scale=0.5] {similar23.png} \end{center}

The problem states that the bases are parallel ($XY \parallel BC$) and also that the subdivision produces \emph{equal areas}.  We are asked to find the ratio $AX:XB$.

\emph{Solution}.

Recall from a previous chapter that for two similar triangles, the altitudes to corresponding sides are in the same ratio as any of the three pairs of sides themselves.  

The altitudes $h$ and $H$ to $BC$ (not drawn) are also in the same ratio as the sides, so let $H = kh$ --- $H$ lies on $BC$.

Then twice the area of the top triangle is $AX \cdot h$ and twice the area of the whole is $AB \cdot H = k AB \cdot h$

We have that the whole is twice the smaller area so the ratio is equal to $2$:
\[ \frac{k AB \cdot h}{AX \cdot h} = 2 \]
But $AB/AX$ is also equal to $k$ so we have that $k^2 = 2$ and $k = \sqrt{2}$.

However, we are not asked for $k$.  Instead, we want the ratio of $AX$ to the smaller piece along the bottom. Let $AX = a$.  Then the whole is $A = ka$.  The difference is the small piece, $XB = A - a = a(k - 1)$.  We need
\[ \frac{a}{a(k - 1)} = \frac{1}{k - 1}  = \frac{1}{\sqrt{2} - 1} \]

\subsection*{problem}

Here's a problem from the web.  Given that $AC \parallel EF$ and $AB \parallel DF$.

We are to prove that the sum of the altitudes of the small triangles is equal to the altitude of the large one.

\begin{center} \includegraphics [scale=0.4] {prob_similar_tri2.png} \end{center}

Informal solution:  The smaller triangles are all similar to $\triangle ABC$ by the alternate interior angles and vertical angle theorems.

For similar triangles, not only are the sides in the same ratio to each other, but so are other measures like the altitudes to a particular side.  So if we label the bases $b_1$ etc., collectively $b_i$ , then we have that
\[ \frac{b_i}{h_i} = \frac{b}{h} \]
\[ b_i = b \cdot h_i/h \]
for each of the $b_i$.

But the sum of the $b_i$ is simply equal to $b$ so
\[ b_1 + b_2 + b_3 = b \cdot (h_1/h + h_2/h + h_3/h) \]
\[ b = b \cdot (h_1/h + h_2/h + h_3/h) \]
\[ 1 = h_1/h + h_2/h + h_3/h \]
\[ h = h_1 + h_2 + h_3 \]

\subsection*{problem}

Given two triangles that are not similar, where the ratio $AD/AB = r$ and $AE/AC = s$.  We are asked to show that 

\[ \frac{\triangle_{ADE}}{\triangle_{ABC}} = rs \] 

\begin{center} \includegraphics [scale=0.4] {similarity_by_area2.png} \end{center}

\emph{Solution}.

Draw $DF \parallel BC$.  Now $\triangle ADF \sim \triangle ABC$ so $AF/AC = r$.

By our previous result
\[ \frac{\triangle_{ADF}}{\triangle_{ABC}} = r^2 \] 

Since $AE/AC = s$ and $AF/AC = r$, $AE/AF = s/r$.  As triangles with a common vertex and bases in that proportion, these areas are in the same proportion:

\[ \frac{\triangle_{ADE}}{\triangle_{ADF}} = \frac{s}{r} \] 

Multiply the two ratios together to obtain

\[ \frac{\triangle_{ADE}}{\triangle_{ABC}} = sr \] 

$\square$

\emph{Solution}.  (alternate).

\begin{center} \includegraphics [scale=0.4] {similarity_by_area2.png} \end{center}

Looking ahead to trigonometry, in any triangle, twice the area can be computed as the product of two sides flanking an angle times the \emph{sine} of the angle.  (The reason is that the altitude to one side, divided by the length of the second side, is defined to be the sine of the angle.)

\[ 2 (ADE) = AD \cdot AE \cdot \sin \angle DAE \]
\[ 2 (ABC) = AB \cdot AC \cdot \sin \angle BAC \]

But $\angle DAE = \angle BAC$, so they have the same sine, and the ratio of areas is just
\[ \frac{AD \cdot AE}{AB \cdot AC} = rs \]

$\square$

To rework this proof in terms of familiar concepts, draw the altitude from $D$ to $AE$, and also the one from $B$ to $AC$

They form similar right triangles including the angle at $A$.  The altitudes scale like $AD/AB = r$.  But the bases scale like $AE/AC = s$.  

And area scales like the product of the two, namely, $rs$.



\end{document}